\section{Descent to the Picard variety}\label{sec:family}

Let $C$ be a smooth projective complex curve of genus $g=b+1\geq1$,
$B=\operatorname{Pic}^{b}(C)$, and $L$ the Poincar\'e line on
$C\times B$, trivialized over a chosen point of $C$.
Use the ordered determinant lines of \cite[Lemma~2.1]{W1} for
the two stable-pair indices. Write $q$ for the original Euler
variable, $t$ for the exchanged grading, $c=\kappa^{-1}$ and
$\sigma^2=\kappa$. For a partition $\lambda$ of $m$, set
$H_\lambda=\sum_{v\in\lambda}h_v$, where $h_v$ and $\ell_v$
denote the hook length and leg length. The scalar removed from each
of its two contributions in passing to the disk calculation is
\begin{equation}\label{eq:family-scalar}
 (-1)^{mb}q^{-bH_\lambda}
 c^{\,b\sum_{v\in\lambda}(\ell_v+1/2)}.
\end{equation}
Here $c^{1/2}=\sigma^{-1}$. Thus disk degree $(n,k)$ corresponds
to $q^{-bH_\lambda+n}t^k$ on the curve.

There is a second normalization in the current calculation. With
$p$ the disk cotangent character and $s=p/c$, the common series of
exchanged degree zero is
\begin{equation}\label{eq:family-axis}
 A_{\lambda;c,s}(x,0)= A_{\lambda;s,c}(0,x)
 =\prod_{v\in\lambda}F_s(x^{h_v}c^{-\ell_v}),\qquad
 F_s(z)=\sum_{i\geq0}\frac{(s;p)_i}{(p;p)_i}z^i .
\end{equation}
Here $(a;p)_i=\prod_{j=0}^{i-1}(1-ap^j)$.
The hook product is \cite[Lemma~5.1]{W1}. Its constant term is one.
Consequently an identity of rational functions in $p,c,x$ for a
coefficient divided by \eqref{eq:family-axis} gives equality of the
full coefficient as a power series in $x$, on multiplying by
\eqref{eq:family-axis}.

\begin{proposition}\label{prop:family}
Fix $d\geq0$ and a partition $\lambda$ of $m$. Suppose the two disk
series $A_{\lambda;c,s}(x,y)$ and $A_{\lambda;s,c}(y,x)$ agree
through $y$-degree $d$ in $\mathbb Q(p,c)[[x,y]]$. Then the
contributions of $\lambda$ to the two family indices agree in
\[
 K^0(B)[\sigma^{\pm1},(1-\sigma^j)^{-1}\ (j\geq1)]
\]
at every $q^{-bH_\lambda+n}t^k$ with $\min(n,k)\leq d$ and
$n,k\geq0$. If the hypothesis holds for every $\lambda\vdash m$,
the two degree-$m$ indices, multiplied by $\sigma^{mb}$, agree
modulo $t^{d+1}$ in
\[
 q^{-bm(m+1)/2}\,K^0(B)[\kappa^{\pm1}][[q,t]].
\]
Here $K^0(B)$ denotes integral algebraic $K$-theory.
\end{proposition}

\begin{proof}
The involution $(x,y,c)\mapsto(y,x,p/c)$ gives the disk identity
also through $x$-degree $d$. For any $(n,k)$ in either strip it
therefore gives every coefficient of the rectangle
$0\leq a\leq n$, $0\leq e\leq k$.
We recall how the two-divisor descent of
\cite[Proposition~7.1]{W1} uses these coefficients.
After removal of \eqref{eq:family-scalar}, the nested divisors of the
fixed boxes and the insertions from the moving factors give classes
on $C^{(n)}\times C^{(k)}\times B$. On disjoint supports their
restrictions are exterior products, with the graded permutation
action on equal pairs of multiplicities.

In an affine coordinate, the derived restriction at a single point
of multiplicities $(a,e)$ is the disk coefficient times
\[
    z^e(p;p)_a(p;p)_e ,
\]
where $z$ is the character of the frame of $L$. The two finite
products are the cotangent Euler classes of the symmetric powers
of the affine line at their origins. Hence the comparison is made
while $p$ and $z$ are independent characters. Expanding the
zero-counting-degree Euler inverses in the volume character, every
coefficient of the restriction is an integral torus character.
It follows that the two restrictions agree integrally. On a cycle
of a permutation of equal support types the supertrace is the
punctual supercharacter evaluated on the product of the torus
entries around that cycle. This includes the cycles of lengths two
and three at the corresponding exchanged degrees.

The derived restrictions detect the two classes on the stack of
coordinate jets and frames used in \cite[\S7.2]{W1}. Indeed,
filter by number of distinct support points. The conormal Euler
of each stratum in its step of the filtration has constant term
one and all other terms of positive simultaneous dilation weight.
It acts injectively on the integral representation group of the
stabilizer, including permutations of equal types. Coherent
localization and the self-intersection formula then give detection
by induction. Pullback to the two symmetric powers uses the actual
cotangent and Poincar\'e lines of the curve. Proper pushforward to
$B$ and restoration of \eqref{eq:family-scalar} give the asserted
diagram equality.

Finally, \cite[Lemma~2.2]{W1} puts each normalized degree sum in
the indicated Laurent polynomial ring coefficientwise. At fixed
$q^j t^k$ with $k\leq d$ a diagram occurs at excess
$n=j+bH_\lambda$ if this is nonnegative. Each such diagram has
been compared. Localization at $1-\sigma^i$, $i>0$, is injective
on $K^0(B)[\sigma^{\pm1}]$, by its extremal unit coefficient,
and $\kappa\mapsto\sigma^2$ is injective there as well. The
summed equality thus holds with the stated integral coefficients.
\end{proof}
